\section{Introduction}
The information available to a planner determines which paths it can evaluate and optimize. Known road topology and edge costs define a fixed optimization problem. A robot with limited vision must also choose where to obtain observations before planning through unknown passages. A locally optimal path may lead into a dead end, requiring return to an observed branching position. Graph search and exploration therefore address distinct, related problems.

This project investigates one C++ implementation of \astar{} in both settings. Road Navigation minimizes distance on a static directed graph. Blind-Alley Robot Navigation plans continuous piecewise-linear motion between exploration targets and returns to observed ancestors using accumulated polygonal observations. Python supplies sensing and exploration around the shared engine. An educational Robot Lab and four-direction blind-alley fixtures remain validated baselines.

Phan et al.~\cite{phan2026} investigate navigation through blind-alley regions and analyze return-path properties under stated geometric assumptions. Their method represents explored space through bundles of line segments and refines paths geometrically. Here, recorded reversible motion supplies a feasible return while \astar{} seeks a shorter verified graph path. This alternative requires explicit sensing and execution assumptions. Comparison must also distinguish computed approximate paths from recorded simulator coordinate transitions.

Accordingly, we examine three questions: how adapters preserve search assumptions; what effort and time geographic guidance saves relative to Dijkstra on the same road graph; and how local robot paths differ at identical observed states and endpoints. The last comparison accounts for representation, validity conventions, and computation stages. Separate online studies examine exploration ordering and executed recovery.

The contribution is a reproducible implementation and controlled evaluation of shared graph search, together with a conditional executed-return certificate. Its scope is established \astar{} theory and observed-graph planning; global unknown-world optimality is not established. Measurements use frozen revision \source{04c4846}.

The report proceeds from theoretical foundations and the problem formulation to system architecture and implementation. The methodology and results distinguish fixed-road shortest paths, shared-state local robot planning, and supporting online behavior. The final sections state limitations, future improvements, and the project conclusions.
